\subsection{A binary composition map above a terminal source block}

\begin{proposition}[Proved: a binary-map deletion or fixed-arrow obstruction]
\label{CP2:A-four-binary-map-dichotomy}
Let $C$ be a category on four objects $p,s,q,r$, in this
order, whose hom-set cardinality matrix is
\[
\begin{pmatrix}
1&0&1&2\\
0&1&1&t\\
0&0&1&2\\
0&0&0&m
\end{pmatrix},\qquad t,m\geq1.
\]
Then either the inclusion
$C\setminus\{p\}\hookrightarrow C$ induces a homotopy
equivalence, or every rational degree-two product on
$BC$ vanishes. The opposite statement also holds.
Consequently no such category with at most eighteen
arrows has classifying-space type $\CP^2$.
\end{proposition}
\begin{proof}
Let $a:p\to q$ be the unique arrow and put
$N=\operatorname{End}_C(r)$. Composition defines an
$N$-equivariant map of two-element left $N$-sets
\[
\lambda:C(q,r)\longrightarrow C(p,r),
\qquad b\longmapsto b\circ a.
\]
Every map between two-element sets is either a bijection
or a constant map.

Suppose first that $\lambda$ is bijective, and let
$j:C\setminus\{p\}\hookrightarrow C$ be the full inclusion.
The comma category $(p\downarrow j)$ has the object
$a:p\to q$ and the two objects $u:p\to r$.
For each $u$ there is exactly one arrow from $a$ to $u$,
because its arrows are precisely the $b:q\to r$ with
$\lambda(b)=u$. The only endomorphism of $a$ is the
identity. Thus $a$ is initial, regardless of the
endomorphisms and other arrows between the objects
over $r$. For an object $c$ of the retained subcategory,
$(c\downarrow j)$ has initial object $\operatorname{id}_c$.
All these comma categories are contractible, so
Quillen's Theorem~A gives the asserted homotopy
equivalence.

Suppose instead that $\lambda$ has constant value $u$.
For every $n\in N$ and $b\in C(q,r)$, associativity gives
\[
n\circ u=n\circ(b\circ a)=(n\circ b)\circ a=u.
\]
Thus $u:p\to r$ is fixed by every left $N$ actor.
The full source block $E=C|_{\{p,s,q\}}$ has terminal
object $q$, so $BE$ is contractible. In the two-block
decomposition of $C$ into $E$ and $N$, the fixed arrow
gives a section of the bridge projection to $BN$.
Corollary~\ref{CP2:A-fixed-arrow}, with
$H^2(BE;\mathbb Q)=0$, therefore makes every degree-two
class relative to the two vertex blocks. Their products
vanish by the suspension quotient in
Lemma~\ref{CP2:A-two-block}.

In the first alternative, deleting $p$ removes its
identity, its single arrow to $q$, and its two arrows
to $r$, hence exactly four arrows. An eighteen-arrow
or smaller projective-plane model would therefore give
one with at most fourteen arrows, contradicting
Theorem~\ref{CP2:A-global-lower-eighteen}. In the second
alternative its nonzero rational square is already
impossible. Opposition preserves the classifying-space
homotopy type and gives the dual statement.
\end{proof}

\begin{corollary}[Proved: the frozen four-object profile $1508$]
\label{CP2:A-four-eighteen-profile-1508}
The frozen eighteen-arrow four-object profile $1508$
does not have classifying-space type $\CP^2$.
\end{corollary}
\begin{proof}
Its hom-set matrix is the displayed matrix with
$m=7$ and $t=2$. Apply
Proposition~\ref{CP2:A-four-binary-map-dichotomy}.
\end{proof}
